\documentclass[11pt]{amsart}
\usepackage[T1]{fontenc}
\usepackage{lmodern}
\usepackage{amsmath,amssymb,mathtools}
\usepackage[margin=1in]{geometry}
\usepackage{microtype}
\usepackage{hyperref}
\hypersetup{colorlinks=true,linkcolor=blue,citecolor=blue,urlcolor=blue,
  pdftitle={Prescribed Goldbach Lengths in the Nonnegative Polynomial Semiring},
  pdfsubject={Additive representations by multiplicative semiring atoms}}
\urlstyle{same}

\newcommand{\N}{\mathbb N}
\newcommand{\Nzero}{\mathbb N_0}
\newcommand{\Z}{\mathbb Z}
\newcommand{\Q}{\mathbb Q}
\newcommand{\Atoms}{\mathcal A}
\newcommand{\Lplus}{\mathsf L_{+}}
\newcommand{\Llaur}{\mathsf L_{\pm}}
\DeclareMathOperator{\supp}{supp}
\newcommand{\coeff}[2]{[x^{#1}]#2}

\newtheorem{theorem}{Theorem}[section]
\newtheorem{lemma}[theorem]{Lemma}
\newtheorem{corollary}[theorem]{Corollary}
\newtheorem{proposition}[theorem]{Proposition}
\theoremstyle{definition}
\newtheorem{example}[theorem]{Example}
\theoremstyle{remark}
\newtheorem{remark}[theorem]{Remark}

\title[Prescribed Goldbach lengths]{Prescribed Goldbach Lengths in the Nonnegative Polynomial Semiring}
\author{Shunyao Yan}
\email{shunyaoyan2009@gmail.com}
\date{Draft revised September 19, 2026}
\subjclass[2020]{Primary 16Y60; Secondary 11P32}
\keywords{Polynomial semiring, multiplicative atom, Goldbach theorem,
additive representation, prescribed length, Laurent polynomial}

\begin{document}

\begin{abstract}
We determine the possible numbers of multiplicative atoms in additive
decompositions of nonconstant polynomials $f\in\Nzero[x]$ with $f(0)\ge2$.
Writing $f=a+bx+P(x)$ with $P\in x^2\Nzero[x]$, we prove that $f$ is a sum
of exactly $k\ge2$ atoms if and only if $k\le a+b$ and
$2k\le a+2b+P(1)$. The proof is constructive and uses elementary parity
criteria. In particular, every such polynomial is a sum of at most two
atoms; the only exceptions to a sum of exactly two are $2+x^r$ with
$r\ge2$. We also record a prescribed-length corollary of Kaplan--Polo's
Laurent two-atom theorem. Atomicity is always in the indicated semiring,
not in $\Q[x]$.
\end{abstract}

\maketitle

\section{Introduction}

Goldbach-type questions ask for representations as sums of multiplicatively
irreducible elements. Hayes~\cite{Hayes} and Pollack~\cite{Pollack} studied
such questions over polynomial rings. In the nonnegative setting,
Liao and Polo~\cite{LiaoPolo} proved that every
$f\in\Nzero[x^{\pm1}]$ with $f(1)>3$ and $|\supp(f)|>1$ is a sum of two
atoms. Kaplan and Polo~\cite{KaplanPolo} recovered this result without
prime-distribution estimates and extended it to Laurent polynomial and
series semidomains. Related work treats group semidomains~\cite{LiMopuriZhang}
and the arithmetic of polynomial semidomains~\cite{GottiPolo}.

In the ordinary semiring $\Nzero[x]$, the only zero-constant atom is $x$,
whereas $x$ is a unit in the Laurent semiring. Thus a zero-constant
polynomial having a term of degree at least two is not a sum of ordinary
polynomial atoms. We instead classify the possible numbers of summands
when the constant coefficient is at least two.

Throughout, $\N=\{1,2,\ldots\}$ and $\Nzero=\{0,1,2,\ldots\}$.
An \emph{atom} is a nonzero nonunit $h$ such that every factorization
$h=uv$ in the indicated semiring has a unit factor. The only unit of
$\Nzero[x]$ is $1$. Write $\Atoms$ for its atoms and define
\[
\Lplus(f)=
\left\{k\in\N: f=h_1+\cdots+h_k\text{ for some }h_1,\ldots,h_k\in\Atoms\right\}.
\]
Repeated summands and prime-constant atoms are permitted. These are
\emph{additive} representation lengths, not multiplicative factorization
lengths. Semiring atomicity is also distinct from rational irreducibility:
$1+x^3$ is a semiring atom, although
$1+x^3=(1+x)(1-x+x^2)$ in $\Q[x]$.

Our main result depends only on the constant coefficient, linear
coefficient, and total mass of the higher-degree coefficients.

\begin{theorem}[Prescribed-length classification]\label{thm:main}
Let $f(x)=a+bx+P(x)\in\Nzero[x]$ be nonconstant, where
$a\ge2$, $b\ge0$, and $P\in x^2\Nzero[x]$. Put $m=P(1)$.
For every integer $k\ge2$,
\begin{equation}\label{eq:capacity}
k\in\Lplus(f)
\quad\Longleftrightarrow\quad
k\le a+b\quad\text{and}\quad 2k\le a+2b+m.
\end{equation}
Equivalently,
\begin{equation}\label{eq:interval}
\begin{gathered}
K(f)=\min\left\{a+b,\left\lfloor\frac{a+2b+m}{2}\right\rfloor\right\},\\
\Lplus(f)\cap\{2,3,\ldots\}=\{2,3,\ldots,K(f)\},
\end{gathered}
\end{equation}
where the right side is empty if $K(f)<2$. Length one occurs precisely
when $f$ itself is an atom.
\end{theorem}

The two inequalities are elementary necessary conditions. Their
sufficiency holds for arbitrary degree and coefficient height, with no
restriction on the positions of the exponents in $P$. The ordinary
classification has a self-contained proof; the Laurent comparison in
Section~\ref{sec:laurent} uses the cited two-atom theorem.

\begin{corollary}[Sharp two-atom statement]\label{cor:two}
Every nonconstant $f\in\Nzero[x]$ with $f(0)\ge2$ is a sum of at most two
atoms. It is a sum of exactly two atoms if and only if
\[
f\ne2+x^r\qquad\text{for every }r\ge2.
\]
The exceptional polynomials $2+x^r$ are themselves atoms.
\end{corollary}

\section{Elementary atom certificates}

For $h=\sum_{j\ge0}h_jx^j\in\Nzero[x]$, define
\[
E(h)=\sum_{j\text{ even}}h_j,\qquad
O(h)=\sum_{j\text{ odd}}h_j.
\]
The constant coefficient contributes to $E(h)$.

\begin{lemma}[Parity certificates]\label{lem:parity}
Suppose $h\in\Nzero[x]$ has positive constant coefficient. Then $h$ is
an atom if either
\[
O(h)=1,\qquad\text{or}\qquad E(h)=1\text{ and }h\ne1.
\]
\end{lemma}

\begin{proof}
If $h=uv$ in $\Nzero[x]$, then $u(0),v(0)>0$, and
\begin{align*}
E(h)&=E(u)E(v)+O(u)O(v),\\
O(h)&=E(u)O(v)+O(u)E(v).
\end{align*}
Both $E(u)$ and $E(v)$ are positive integers. If $O(h)=1$, exactly one
term in the second identity is one. After interchanging $u,v$, we have
$O(u)=0$ and $E(u)=1$. Since $u(0)>0$, this forces $u=1$.
If $E(h)=1$, the first identity forces $E(u)=E(v)=1$ and
$O(u)O(v)=0$. A factor with no odd terms has exactly one even coefficient
unit, necessarily at degree zero, and is therefore $1$.
In either case $h$ is a nonunit, so it is an atom.
\end{proof}

\begin{lemma}[Exponent rescaling]\label{lem:rescale}
If $H\in\Nzero[x]$ has positive constant coefficient and $D\ge1$, then
$H(x^D)$ is an atom if and only if $H(x)$ is an atom. In particular,
$c+x^r$ is an atom for all integers $c,r\ge1$.
\end{lemma}

\begin{proof}
Suppose $H(x^D)=uv$ in $\Nzero[x]$. Both factors have positive constant
coefficient. Every exponent in either factor's support also occurs in
the product: multiply that term by the other factor's nonzero constant
term and use nonnegativity. Hence every factor exponent is divisible by
$D$, so $u=U(x^D)$ and $v=V(x^D)$ for some $U,V\in\Nzero[x]$.
This factorization descends to $H=UV$, and units are preserved by
substitution. The converse follows by substituting in a factorization of
$H$. Finally, $c+x$ satisfies Lemma~\ref{lem:parity}; rescaling gives
$c+x^r$.
\end{proof}

\begin{lemma}[Two-atom construction]\label{lem:two}
Let $A\ge2$, let $0\ne R\in x\Nzero[x]$, and suppose $A+R(1)\ge4$.
Then $A+R$ is a sum of two atoms of $\Nzero[x]$.
\end{lemma}

\begin{proof}
If $R(1)=1$, then $R=x^r$ and $A\ge3$. The decomposition
\[
A+x^r=2+\bigl((A-2)+x^r\bigr)
\]
uses the prime constant $2$ and an atom from Lemma~\ref{lem:rescale}.

Suppose $R(1)\ge2$. Let $D$ be the largest power of two dividing every
exponent in $\supp(R)$. Write $R(x)=T(x^D)$. At least one exponent of
$T$ is odd; otherwise $2D$ would divide every exponent of $R$.
Decompose
\[
T=T_{\mathrm e}+T_{\mathrm o},
\]
where $T_{\mathrm e}$ has positive even exponents and $T_{\mathrm o}$ has
odd exponents. Thus $T_{\mathrm o}(1)\ge1$.

If $T_{\mathrm o}(1)\ge2$, choose an odd exponent $r$ occurring in
$T_{\mathrm o}$. Then
\begin{equation}\label{eq:split-odd}
A+T=
\bigl((A-1)+x^r+T_{\mathrm e}\bigr)
+\bigl(1+T_{\mathrm o}-x^r\bigr).
\end{equation}
The first summand has odd coefficient sum one. The second has even
coefficient sum one and is not the unit, because
$T_{\mathrm o}(1)\ge2$. Both have positive constant coefficient and are
atoms by Lemma~\ref{lem:parity}.

If $T_{\mathrm o}(1)=1$, write $T_{\mathrm o}=x^r$. Then
$T_{\mathrm e}\ne0$, since $T(1)\ge2$. Choose an even exponent $s$
occurring in $T_{\mathrm e}$. We have
\begin{equation}\label{eq:split-even}
A+T=
\bigl((A-1)+x^s\bigr)
+\bigl(1+x^r+T_{\mathrm e}-x^s\bigr).
\end{equation}
The first summand is an atom by Lemma~\ref{lem:rescale}; the second has
odd coefficient sum one and is an atom by Lemma~\ref{lem:parity}.
Every coefficient in these decompositions is nonnegative. Substituting
$x^D$ preserves atomicity by Lemma~\ref{lem:rescale} and gives the
required decomposition of $A+R$.
\end{proof}

\section{Proof of the prescribed-length classification}

For $h=\sum_{j\ge0}h_jx^j$, introduce the additive weights
\begin{equation}\label{eq:weights}
U(h)=h_0+h_1,\qquad
W(h)=h_0+2h_1+\sum_{j\ge2}h_j=h(1)+h_1.
\end{equation}

\begin{lemma}[Necessary capacity inequalities]\label{lem:capacity}
Every atom $h\in\Nzero[x]$ satisfies $U(h)\ge1$ and $W(h)\ge2$.
Consequently, every sum of $k$ atoms satisfies $U(f)\ge k$ and
$W(f)\ge2k$.
\end{lemma}

\begin{proof}
The only zero-constant atom is $x$: a nonzero polynomial with zero
constant coefficient factors as $xq$, and atomicity forces $q=1$.
The constant atoms are exactly the positive integer primes. Any other
atom has positive constant coefficient and at least one positive-degree
term. For $x$ the weights are $(1,2)$; for a prime constant they are
both at least two; and for a positive-constant nonconstant atom the
inequalities are immediate. Additivity proves the assertion for a sum.
\end{proof}

\begin{proof}[Proof of Theorem~\ref{thm:main}]
Necessity follows from Lemma~\ref{lem:capacity}, since $U(f)=a+b$ and
$W(f)=a+2b+m$.

For sufficiency, fix an integer $k\ge2$ satisfying the two inequalities,
and induct on $k$. If $k=2$ and $f(1)\ge4$, apply Lemma~\ref{lem:two}.
Otherwise $f(1)=3$, so $f=2+x^r$. The condition $W(f)\ge4$ forces
$r=1$, and $2+x$ is the sum of the two atoms $2$ and $x$.

Now let $k\ge3$. We remove one atom so that the remainder is
nonconstant, has constant coefficient at least two, and satisfies the
capacity inequalities for $k-1$.

\smallskip
\noindent\textit{Case 1: $a\ge k+1$.}
Remove $2$. The remainder has constant coefficient
$a-2\ge k-1\ge2$ and stays nonconstant. Moreover,
\[
U(f-2)=a+b-2\ge k-1,\qquad
W(f-2)=W(f)-2\ge2(k-1).
\]

\smallskip
\noindent\textit{Case 2: $a\le k$ and $b\ge1$.}
Remove $x$. The constant coefficient is unchanged, and
\[
U(f-x)=U(f)-1\ge k-1,\qquad
W(f-x)=W(f)-2\ge2(k-1).
\]
If $f-x$ were constant, then $f=a+x$, giving
$2k\le a+2\le k+2$, impossible for $k\ge3$. Thus the remainder is
nonconstant.

\smallskip
\noindent\textit{Case 3: $a\le k$ and $b=0$.}
The first inequality gives $a=k$, and the second gives $m\ge k$.
Choose $r\ge2$ with $\coeff rP>0$, and remove the atom $1+x^r$.
The remainder has constant coefficient $k-1\ge2$ and higher-degree
mass $m-1\ge k-1>0$. Its weights satisfy
\[
U\bigl(f-(1+x^r)\bigr)=k-1,\qquad
W\bigl(f-(1+x^r)\bigr)=W(f)-2\ge2(k-1).
\]
In all three cases, the induction hypothesis writes the remainder as a
sum of $k-1$ atoms. Adjoining the removed atom proves sufficiency.
Formula~\eqref{eq:interval} follows directly from~\eqref{eq:capacity},
and the assertion about length one follows from the definition.
\end{proof}

The proof is constructive: repeatedly remove $2$, $x$, or $1+x^r$
according to the three cases, and finish with the explicit two-atom
decomposition. No search for integer prime pairs is required.

\section{Sharp consequences in the ordinary semiring}

\begin{proof}[Proof of Corollary~\ref{cor:two}]
For $k=2$, the bound $2\le a+b$ is automatic. For a nonconstant target
with $a\ge2$, the other inequality fails exactly when
$a=2$, $b=0$, and $m=1$, that is, when $f=2+x^r$ with $r\ge2$.
These exceptional polynomials are atoms by Lemma~\ref{lem:rescale}.
\end{proof}

\begin{corollary}[Constant coefficient two]\label{cor:constant-two}
Let $f=2+bx+P$ be nonconstant, where $b\ge0$ and $P\in x^2\Nzero[x]$.
Put $m=P(1)$. Then $\Lplus(f)\cap\{2,3,\ldots\}=\{2,\ldots,K(f)\}$,
where
\[
K(f)=
\begin{cases}
b+1,&m=0\text{ or }1,\\
b+2,&m\ge2.
\end{cases}
\]
The interval is empty when $K(f)=1$.
\end{corollary}

\begin{proof}
Theorem~\ref{thm:main} gives
\[
K(f)=\min\left\{b+2,\ b+1+\left\lfloor\frac m2\right\rfloor\right\}.
\qedhere
\]
\end{proof}

\begin{example}[Why the linear coefficient matters]\label{ex:linear}
The polynomials $2+2x$ and $2+2x^r$ for $r\ge2$ have the same total
coefficient sum, but
\[
\Lplus(2+2x)=\{2,3\},\qquad
\Lplus(2+2x^r)=\{2\}.
\]
For the first, $2+x+x$ is a three-atom representation. For the second,
$U(f)=2$ restricts every representation to at most two summands.
Neither polynomial is an atom, since both have the nonunit constant
factor $2$.
\end{example}

\begin{remark}[The constant-one boundary]\label{rem:constant-one}
The hypothesis $f(0)\ge2$ cannot be omitted. For instance,
\[
\Lplus(1+7x+6x^2)=\{2,4,5,6,7,8\},
\]
which has a gap. In any atomic representation of a constant-one
polynomial, exactly one summand has positive constant coefficient; all
others must be $x$. Thus a length-$k$ representation of this example
exists exactly when
\[
1+(8-k)x+6x^2
\]
is an atom, with $1\le k\le8$.
For a nonnegative integer $t$, the polynomial $1+tx+6x^2$ is reducible in
$\Nzero[x]$ precisely when
\[
1+tx+6x^2=(1+ux)(1+vx),\qquad uv=6,
\]
so $t=u+v\in\{5,7\}$. These exclude exactly $k=3,1$, proving the
claimed length set. The constant-one case therefore has a different
character from Theorem~\ref{thm:main}.
\end{remark}

\section{A prescribed-length consequence of Laurent Goldbach theorems}
\label{sec:laurent}

Let
\[
R_t=\Nzero[x_1^{\pm1},\ldots,x_t^{\pm1}],\qquad t\ge1.
\]
Its units are precisely the coefficient-one monomials. A Laurent
polynomial is a finite sum with exponent vectors in $\Z^t$; write
$\nu(F)=F(1,\ldots,1)$ for its coefficient mass. Define $\Llaur(F)$
to be the set of positive integers $k$ for which $F$ is a sum of exactly
$k$ multiplicative atoms of $R_t$.

Liao--Polo's univariate theorem~\cite{LiaoPolo}, and its
multivariate extension by Kaplan and Polo~\cite[Corollary~3.8]{KaplanPolo},
state that every $F\in R_t$ with at least two support exponents and
$\nu(F)>3$ is a sum of two atoms. The following is an elementary
consequence, not a separate two-atom theorem.

\begin{proposition}[Laurent prescribed lengths]\label{prop:laurent}
Let $F\in R_t$ have at least two distinct support exponents. For every
integer $k\ge2$,
\[
k\in\Llaur(F)\quad\Longleftrightarrow\quad2k\le\nu(F).
\]
Length one occurs precisely when $F$ itself is a Laurent atom.
\end{proposition}

\begin{proof}
Every Laurent atom has mass at least two, since every mass-one element
is a unit. This proves necessity.

Suppose $\nu(F)\ge2k$. Reserve one coefficient unit at each of two
distinct support exponents. From the remaining units, select
$2(k-2)$ and partition them into $k-2$ pairs. Each pair is an atom:
it has mass two, whereas a product of two nonzero nonunits has mass
at least four.

After removing these atoms, the remainder $G$ has
\[
\nu(G)=\nu(F)-2(k-2)\ge4,
\]
and retains two support exponents because the reserved units were
not removed. Kaplan--Polo's theorem expresses $G$ as a sum of two
Laurent atoms. Adjoining the removed $k-2$ atoms gives exactly $k$
summands.
\end{proof}

\begin{remark}[Why localization changes the answer]
For nonmonomial Laurent polynomials, the admissible representation
lengths are characterized entirely by the coefficient mass through
the inequality $2k\le\nu(F)$. In contrast, the ordinary polynomial
semiring $\Nzero[x]$ exhibits a richer structure. Every atom other
than $x$ contributes at least one unit to the constant coefficient,
while each occurrence of $x$ contributes to the linear coefficient.
This introduces the additional sharp constraint $k\le a+b$.
Furthermore, since $x$ is an atom of mass one in $\Nzero[x]$
but becomes a unit in the Laurent semiring, the corresponding
weighted constraint in the ordinary setting is
$2k\le f(1)+b$. Together, Theorem~\ref{thm:main} and
Proposition~\ref{prop:laurent} provide a complete comparison
of prescribed representation lengths in these two settings,
identifying precisely how localization changes the additive
structure of multiplicative atoms.
\end{remark}

\section*{Acknowledgment}
AI tools like GPT-5.5, GPT-5.6, GPT-6 assisted mathematical exploration, proof
checking, and paper preparation.

\begin{thebibliography}{9}

\bibitem{GottiPolo}
F. Gotti and H. Polo,
\emph{On the arithmetic of polynomial semidomains},
Forum Math. \textbf{35} (2023), no.~5, 1179--1197.
\href{https://doi.org/10.1515/forum-2022-0091}{doi:10.1515/forum-2022-0091}.

\bibitem{Hayes}
D.~R. Hayes,
\emph{A Goldbach theorem for polynomials with integral coefficients},
Amer. Math. Monthly \textbf{72} (1965), no.~1, 45--46.
\href{https://doi.org/10.2307/2312999}{doi:10.2307/2312999}.

\bibitem{KaplanPolo}
N. Kaplan and H. Polo,
\emph{A Goldbach theorem for Laurent series semidomains},
New York J. Math. \textbf{31} (2025), 1409--1426.
\url{https://nyjm.albany.edu/j/2025/31-53v.pdf}.

\bibitem{LiMopuriZhang}
E. Li, A. Mopuri, and C. Zhang,
\emph{Goldbach theorems for group semidomains},
preprint, arXiv:2412.00590 (2024).
\url{https://arxiv.org/abs/2412.00590}.

\bibitem{LiaoPolo}
S. Liao and H. Polo,
\emph{A Goldbach theorem for Laurent polynomials with positive integer coefficients},
Amer. Math. Monthly \textbf{131} (2024), no.~8, 704--711.
\href{https://doi.org/10.1080/00029890.2024.2363725}{doi:10.1080/00029890.2024.2363725}.

\bibitem{Pollack}
P. Pollack,
\emph{On polynomial rings with a Goldbach property},
Amer. Math. Monthly \textbf{118} (2011), no.~1, 71--77.
\href{https://doi.org/10.4169/amer.math.monthly.118.01.071}{doi:10.4169/amer.math.monthly.118.01.071}.

\end{thebibliography}

\end{document}
